
\section{Introduction}

\subsection{Motivation}

Multi-head self-attention is the core operation of every modern transformer, 
but it is also the largest consumer of parameters and compute. A standard 
12-layer BERT-base model has 12 attention heads per layer, each with its own 
$W_Q, W_K, W_V, W_O$ projections totaling 4.7M parameters per layer.

\textbf{Head redundancy is well documented.} Voita et al. (2019) find that over 
half of the heads in BERT can be pruned without significant accuracy loss. 
Michel et al. (2019) show that in many layers, a \textbf{single head} is 
sufficient. Clark et al. (2019) characterize attention patterns as 
"positional", "syntactic", or "rare-word" -- with the majority falling into 
redundant categories.

\textbf{The naive approach -- prune heads -- loses information.} Simply removing 
a head discards whatever information it carried. If the head was doing 
something useful, accuracy drops.

\textbf{The alternative -- fine-tune the full model -- is expensive.} Retraining a 
pruned model requires substantial compute and full-parameter optimization.

\textbf{The opportunity.} What if we could \textbf{wrap} the pretrained attention with 
a small trainable adapter, and let the adapter learn to compensate for 
compressed head capacity? This is the idea behind \textbf{Squeezing Heads}.

\subsection{Problem Statement}

Given a pretrained transformer with $H$ attention heads per layer, we seek 
a modified module $\mathcal{A}'$ such that:

\begin{enumerate}
\item \textbf{Compressible:} Effective head dimension can be reduced by a factor $s$.
\item \textbf{Distillable:} Trainable adapters can recover lost capacity.
\item \textbf{Drop-in:} Pretrained weights are copied exactly; no architecture change.
\item \textbf{Portable:} Works across flat, nested, and causal architectures.
\item \textbf{Quality:} $\Delta\text{Acc} \leq 5\%$ on SST-2.

\end{enumerate}
\subsection{Our Approach: Squeezing Heads}

\textbf{Squeezing Heads} replaces each attention projection with a 
\textbf{LoRA-wrapped linear layer}:

\[
W_\text{eff} = W_\text{frozen} + \frac{\alpha}{r} B A
\]

where:
\begin{itemize}
\item $W_\text{frozen}$ is the pretrained projection (unchanged).
\item $B A$ is a low-rank adapter with rank $r$ (trainable).
\item $\alpha$ is the LoRA scaling factor.

\end{itemize}
\textbf{Figure 1.1} compares three head-compression strategies:

\begin{table}[H]
\centering
\small
\begin{tabular}{llll}
\toprule
\textbf{Strategy} & \textbf{Prunes heads?} & \textbf{Trains full model?} & \textbf{Recovers info?} \\
\midrule
Head pruning \& Yes \& No &  \\
Full fine-tune \& Yes \& Yes & (expensive) \\
\textbf{Squeezing Heads} \& Optional & \textbf{LoRA only} &  \\
\bottomrule
\end{tabular}
\end{table}

\textbf{Key design choices:}

\begin{enumerate}
\item \textbf{Wrapper, not replacement.} The pretrained attention is preserved
  exactly. Only the LoRA parameters differ from initialization.
\item \textbf{Per-projection LoRA.} Each of $W_Q, W_K, W_V, W_O$ has its own LoRA
  adapter, allowing fine-grained capacity allocation.
\item \textbf{GQA-aware pooling.} When the source model uses GQA (Qwen, LLaMA),
  we pool the source KV heads into the target head count.
\item \textbf{Principal head preservation.} When the head dimension is reduced, we
  preserve the leading head dimensions (the principal components) rather
  than random slices.

\end{enumerate}
\subsection{Contributions}

\begin{itemize}
\item \textbf{A LoRA-wrapped attention for head distillation.} Squeezing Heads
  preserves pretrained weights exactly while enabling low-rank adaptation.
\item \textbf{A GQA-aware copy procedure.} We handle the Grouped Query Attention
  case by pooling source KV heads.
\item \textbf{Multi-architecture support.} DistilBERT, RoBERTa, GPT-2, and Qwen with
  RoPE, GQA, and causal masking preserved.
\item \textbf{A systematic evaluation.} Benchmarked on DistilBERT, RoBERTa, and GPT-2
  on SST-2 with ablations over LoRA rank, head-dim reduction, and GQA
  pooling.
\item \textbf{Full reproducibility.} Code, checkpoints, and evaluation scripts
  released under MIT.

\end{itemize}
\subsection{Paper Organization}

Section 2 reviews related work. Section 3 provides background on attention, 
LoRA, and head pruning. Section 4 presents Squeezing Heads. Section 5 
describes integration with four architectures. Section 6 details the 
experimental setup. Section 7 reports results. Section 8 analyzes the 
contribution of LoRA and head-dim reduction. Section 9 discusses practical 
guidance. Section 10 lists limitations. Section 11 concludes.



\subsection{Notation}

\begin{lstlisting}

---
\end{lstlisting}